\documentclass[UTF8,a4paper,11pt]{ctexart}

\usepackage[margin=2.35cm]{geometry}
\usepackage{amsmath,amssymb,mathtools}
\usepackage{booktabs,longtable,array}
\usepackage{enumitem}
\usepackage{xcolor}
\usepackage[hidelinks]{hyperref}

\newcommand{\SpecOp}{\mathsf{S}}
\newcommand{\ProofOp}{\mathsf{P}}
\newcommand{\FindBug}{\mathsf{F}}
\newcommand{\Pre}{\mathsf{Pre}}
\newcommand{\Post}{\mathsf{Post}}
\newcommand{\Body}{\mathsf{Body}}
\newcommand{\Anc}{\mathsf{Anc}}
\newcommand{\Desc}{\mathsf{Desc}}
\newcommand{\SCC}{\mathsf{SCC}}
\newcommand{\Assume}{\mathsf{Assume}}
\newcommand{\Fix}{\mathsf{Fix}}

\title{调用图上的规约生成与证明生成三条宏观路线\\
\large 双访问序列、Bug 传播与最坏工程量}
\author{QCIP 技术调研}
\date{2026 年 9 月}

\begin{document}
\maketitle

\begin{abstract}
本文只研究 annotation 驱动验证在调用图上的宏观调度问题。调用图节点统一表示函数或循环；每个节点有两次主访问：规约访问 $\SpecOp(v)$ 和证明访问 $\ProofOp(v)$。本文形式化三条路线：（I）自顶向下、先在假设 callee 合同成立时完成 caller 证明，再递归证明 callee；（II）自底向上、先生成并证明 callee，再生成并证明 caller；（III）先自顶向下生成全图规约，再自底向上完成证明。顶层规约既可能外部给定，也可能由路线生成。本文还定义作用于节点的 $\FindBug(v)$ 操作：spec bug 沿反向调用边向上传播，proof bug 沿正向调用边向下定位，直到传播条件消失或到达叶节点。最后给出链、分叉、菱形共享、递归强连通分量和包含循环的调用图序列，并分析节点化实现与按路径重复实现的最坏工程量。
\end{abstract}

\section{统一模型}

\subsection{扩展调用图}

令扩展调用图为 $G=(V,E)$，其中 $V=V_F\uplus V_L$：$V_F$ 是函数节点，$V_L$ 是循环节点。边 $(u,v)\in E$ 表示 $u$ 的验证依赖 $v$ 的合同：
\begin{itemize}
  \item 若 $u,v\in V_F$，它是普通函数调用；
  \item 若 $v\in V_L$，它表示函数或外层循环进入该循环；
  \item 嵌套循环用 $L_1\to L_2$ 表示；
  \item 循环不再被单独抽象成特殊阶段，而是与 callee 一样拥有规约和证明操作。
\end{itemize}

循环节点 $L$ 的规约就是 invariant、frame、退出关系等组成的合同 $C_L$；其证明 $\ProofOp(L)$ 包含初始化、保持性和退出三个义务。因此后续所有定义同时适用于函数和循环。

\subsection{每个节点的两次主访问}

每个节点 $v\in V$ 有两个版本化操作：
\[
  \SpecOp(v):\quad \text{生成、接收或固定合同 }C_v,
\]
\[
  \ProofOp(v):\quad
  \Assume\bigl(\{C_w\mid(v,w)\in E\}\bigr)
  \vdash \{\Pre_v\}\Body_v\{\Post_v\}.
\]
这里“证明 $v$”允许模块化地假设直接 callee 的合同成立；但整个验证完成前，每个被依赖合同都必须最终被证明或被明确标记为可信外部边界。Dafny 一类模块化验证器正是逐方法验证，并只使用其他方法的规格\cite{Leino2010Dafny}。

对 root $r$，有两种输入模式：
\begin{enumerate}[label=(\alph*)]
  \item \textbf{root spec 给定}：$\SpecOp(r)=\Fix(C_r^{\mathrm{user}})$，不计算新语义；
  \item \textbf{root spec 未给定}：路线必须生成 $C_r$。此时仍需外部意图、可观察行为模板或“推断实现摘要”的策略，否则正确的功能规约并不唯一。
\end{enumerate}

\subsection{有效完成条件}

设 $T$ 是包含 $\SpecOp(v)$、$\ProofOp(v)$ 和可选 $\FindBug(v)$ 的事件序列。一次 run 完成需要：
\[
  \forall v\in V_{\mathrm{reachable}}.\quad
  \SpecOp(v)\text{ 已固定}\land \ProofOp(v)\text{ 已通过},
\]
并且每个证明绑定的 callee spec version 与最终固定版本一致。若 $C_v$ 变化，则所有依赖旧版本的证明不能继续作为当前结论。

\section{三条宏观路线}

\subsection{路线 I：自顶向下，caller proof 先于 callee proof}

该路线先确定当前节点规约，然后在假设直接 callee 合同成立的情况下完整证明当前节点；当前节点证明完成后，才递归证明刚才假设的 callee 合同。

对节点 $v$ 的递归过程定义为：
\begin{enumerate}[label=I\arabic*.]
  \item 若 $C_v$ 未给定或未固定，执行 $\SpecOp(v)$；
  \item 分析 $\Body_v$ 的调用点，为每个直接 callee $w$ 生成或固定 $C_w$；
  \item 在假设所有 $C_w$ 成立的条件下，完成 $\ProofOp(v)$；
  \item 按选定次序对尚未证明的 $w$ 递归执行本过程；
  \item 所有 callee 均通过后，$\ProofOp(v)$ 才从“条件成立”升级为全局 accepted。
\end{enumerate}

抽象序列为
\[
  \SpecOp(v),
  \SpecOp(w_1),\ldots,\SpecOp(w_k),
  \ProofOp(v\mid C_{w_1},\ldots,C_{w_k}),
  T_I(w_1),\ldots,T_I(w_k).
\]

这条路线允许先验证顶层设计是否在\emph{假定接口}下闭合，适合 caller 驱动的接口设计。但 callee 证明较晚：一旦 $C_w$ 不可实现或必须修改，已经完成的 caller proof 可能全部失效。模块化逻辑允许这种证明顺序；工程风险来自合同版本失效，而不是逻辑规则本身。

\paragraph{root spec 两种情况。}
若给定 $C_r$，第一步只是登记并冻结用户规格；若未给定，则先从 root 需求生成 $C_r$。其余遍历相同。

\subsection{路线 II：自底向上，callee spec/proof 先于 caller}

该路线先递归处理所有 callee；当前节点只能使用已经证明的 callee 合同来生成自己的合同并完成证明。

\begin{enumerate}[label=II\arabic*.]
  \item 对每个直接 callee $w$，先递归完成 $T_{II}(w)$；
  \item 基于 $\Body_v$、已证明的 $C_w$ 和可选外部需求生成或确认 $C_v$；
  \item 使用已证明的 callee 合同完成 $\ProofOp(v)$。
\end{enumerate}

抽象序列为
\[
  T_{II}(w_1),\ldots,T_{II}(w_k),
  \SpecOp(v\mid C_{w_1},\ldots,C_{w_k}),
  \ProofOp(v).
\]

该路线与组合式 bottom-up summary analysis 最接近；bi-abduction 工作会独立分析过程并组合被调过程摘要\cite{Calcagno2009BiAbduction}。

\paragraph{root spec 未给定。}
从叶节点开始根据实现生成并证明摘要，逐层上推，最后得到并证明 $C_r$。

\paragraph{root spec 给定。}
$C_r^{\mathrm{user}}$ 在 run 开始时登记为最终约束，但不用于向下派生 callee 合同；先完成底层合同，最后检查实现摘要是否足以满足 $C_r^{\mathrm{user}}$。其优点是底层结论先稳定，缺点是顶层不匹配可能很晚才暴露。

\subsection{路线 III：spec 自顶向下，proof 自底向上}

该路线明确分成两次全图遍历。

\paragraph{第一遍：Spec pass。}
从 root 出发，根据 caller 规约和调用点需求生成 callee 规约：
\[
  \SpecOp(r)\prec \SpecOp(v)
  \quad\text{for every reachable }r\leadsto v.
\]
对共享 callee，必须先收集多个 caller 的需求，再固定一次稳定合同；不能因第一次到达就忽略后续入边。

\paragraph{第二遍：Proof pass。}
按逆依赖方向证明：
\[
  (v,w)\in E\Longrightarrow \ProofOp(w)\prec\ProofOp(v).
\]
所以整体是
\[
  \underbrace{\SpecOp\text{ 的拓扑序}}_{\text{需求向下}}
  \quad;\quad
  \underbrace{\ProofOp\text{ 的逆拓扑序}}_{\text{证明向上}}.
\]
给定模块接口后推断内部过程合同和循环 annotation 的 intra-module inference 与该路线的 spec pass 思想相近\cite{Lahiri2009IntraModule}；证明阶段仍采用模块化 callee-first。

该路线在进入任何 caller proof 前先固定可达规约，并在 caller proof 前先验证 callee，因此减少“先完成 caller proof、后发现 callee 合同不可实现”的重做。但它可能在尚无证明反馈时生成较多规约。

\section{典型调用图的遍历序列}

以下忽略同层节点之间无依赖时的可交换顺序。$S_X$ 表示 $\SpecOp(X)$，$P_X$ 表示 $\ProofOp(X)$。若 root spec 给定，把 $S_A$ 理解为 $\Fix(C_A^{\mathrm{user}})$；若未给定，则理解为生成 $C_A$。

\subsection{链：$A\to B\to C$}

\begin{center}
\begin{tabular}{ll}
\toprule
路线 & 一个合法序列\\
\midrule
I & $S_A,S_B,P_A,S_C,P_B,P_C$\\
II & $S_C,P_C,S_B,P_B,S_A,P_A$\\
III & $S_A,S_B,S_C,P_C,P_B,P_A$\\
\bottomrule
\end{tabular}
\end{center}

路线 I 中，$P_A$ 假设 $C_B$；证明 $B$ 时才生成 $C_C$。路线 II 的每个节点在后序遍历中连续执行 spec/proof。路线 III 是严格的前序 spec、后序 proof。

\subsection{分叉：$A\to B$ 且 $A\to C$}

\begin{center}
\begin{tabular}{ll}
\toprule
路线 & 一个合法序列\\
\midrule
I & $S_A,S_B,S_C,P_A,P_B,P_C$\\
II & $S_B,P_B,S_C,P_C,S_A,P_A$\\
III & $S_A,S_B,S_C,P_B,P_C,P_A$\\
\bottomrule
\end{tabular}
\end{center}
其中 $B,C$ 可并行；路线 I 只有在 $P_B,P_C$ 都通过后，$P_A$ 才最终解除假设。

\subsection{菱形共享：$A\to B,C$，$B\to D$，$C\to D$}

\begin{center}
\begin{tabular}{ll}
\toprule
路线 & 一个合法序列\\
\midrule
I & $S_A,S_B,S_C,P_A,S_D,P_B,P_D,P_C$\\
II & $S_D,P_D,S_B,P_B,S_C,P_C,S_A,P_A$\\
III & $S_A,S_B,S_C,S_D,P_D,P_B,P_C,P_A$\\
\bottomrule
\end{tabular}
\end{center}

该图揭示三条路线的关键差异：
\begin{itemize}
  \item 路线 I 可能在证明 $B$ 时根据第一个调用上下文固定 $C_D$；之后证明 $C$ 才发现第二个需求，因而必须加强或增加 $D$ 的派生规约。
  \item 路线 II 先从 $D$ 实现得到稳定能力，再供 $B,C$ 使用，但这个能力可能不够满足顶层意图。
  \item 路线 III 在 spec pass 中汇总 $B\to D$ 与 $C\to D$ 两条入边的需求后只固定一次 $C_D$，随后从 $D$ 向上证明。
\end{itemize}

\subsection{递归：$A\leftrightarrow B$}

有环图不存在普通拓扑序。先将 $G$ 缩点为 SCC DAG。对于 $K=\{A,B\}=\SCC(A)$，必须联合固定合同族 $C_K=\{C_A,C_B\}$，并证明
\[
  \Assume(C_A,C_B)\vdash
  \{\Pre_A\}\Body_A\{\Post_A\}
  \land
  \{\Pre_B\}\Body_B\{\Post_B\}.
\]
三条路线在 SCC 内退化为：
\begin{center}
\begin{tabular}{ll}
\toprule
路线 & SCC 内处理\\
\midrule
I & $S_A,S_B,P_A,P_B$，两个 proof 完成后联合提交\\
II & 无真正叶节点；必须先求 $C_A,C_B$ 的联合不动点，再联合证明\\
III & $S_A,S_B$ 联合冻结，然后 $P_A,P_B$ 联合提交\\
\bottomrule
\end{tabular}
\end{center}
SCC 之间仍可分别应用三条宏观路线。也就是说，“任意 CG 的遍历序列”应定义在缩点 DAG 上，SCC 内使用联合事件块。

\subsection{循环作为 callee}

设函数 $A$ 包含外层循环 $L_1$，$L_1$ 调用函数 $B$ 并包含内层循环 $L_2$：
\[
  A\to L_1,\qquad L_1\to B,\qquad L_1\to L_2.
\]
则合法序列为：
\begin{center}
\begin{tabular}{ll}
\toprule
路线 & 一个合法序列\\
\midrule
I & $S_A,S_{L_1},P_A,S_B,S_{L_2},P_{L_1},P_B,P_{L_2}$\\
II & $S_B,P_B,S_{L_2},P_{L_2},S_{L_1},P_{L_1},S_A,P_A$\\
III & $S_A,S_{L_1},S_B,S_{L_2},P_B,P_{L_2},P_{L_1},P_A$\\
\bottomrule
\end{tabular}
\end{center}
这里 $S_{L_i}$ 就是生成 invariant/loop contract，$P_{L_i}$ 就是证明初始化、保持和退出，不需要第四条“loop 专用路线”。

\section{FindBug 操作与传播方向}

\subsection{节点化 Bug 操作}

$\FindBug(v)$ 作用于一个节点及其当前 spec/proof version。返回值至少区分：
\[
  \FindBug(v)\in
  \{\mathsf{no\ bug},\mathsf{spec\ bug},\mathsf{proof\ bug},
    \mathsf{implementation\ bug},\mathsf{unknown}\}.
\]
本文按用户指定的两类传播规则建模。

\subsection{Spec bug 向上传播}

若 $C_v$ 错误、过弱、过强或与调用点需求不一致，则受影响集合位于反向可达闭包：
\[
  U(v)=\{u\mid u\leadsto v\}\cup\{v\}=\Anc(v)\cup\{v\}.
\]
传播不是无条件覆盖所有祖先。定义依赖谓词
\[
  \mathsf{Uses}(u,v,\delta C_v),
\]
表示 $u$ 的 spec/proof 实际使用了本次变化的合同分量。只沿满足该谓词的反向边传播：
\[
  (u,v)\in E\land\mathsf{Uses}(u,v,\delta C_v)
  \Longrightarrow \mathsf{stale}(u).
\]
当某个 caller 不使用变化部分、存在不变的派生规约适配层，或到达 root 时停止。

\subsection{Proof bug 向下传播}

若 $\ProofOp(v)$ 失败且当前尚不能断定是 $v$ 本地实现错误，则沿 callee/loop 边向下查找：
\[
  D(v)=\{w\mid v\leadsto w\}\cup\{v\}=\Desc(v)\cup\{v\}.
\]
定义 $\mathsf{Delegates}(v,w,o)$ 表示未闭合义务 $o$ 依赖 $w$ 的合同或证明。仅当该谓词成立时继续传播到 $w$。停止条件为：
\begin{enumerate}
  \item 在当前节点已定位为本地代码/证明脚本错误；
  \item callee 的 accepted proof 与当前版本匹配，且义务不依赖其未证明性质；
  \item 到达叶节点；
  \item 发现真正问题是 spec bug，随后切换为向上传播。
\end{enumerate}

\subsection{版本失效原则}

每个结果绑定
\[
  \mathsf{version}(v)=
  H(C_v,\Body_v,\{\mathsf{version}(w):(v,w)\in E\}).
\]
spec bug 修改 $C_v$ 后，旧 $\ProofOp(v)$ 和所有使用旧版本的 caller proof 自动 stale。节点版本和边依赖必须去重；否则同一个共享 callee 会沿多条路径被重复传播。

\section{工程量模型}

\subsection{无返工时的基线}

令 $n=|V|$、$m=|E|$，并令 $s_v,p_v$ 分别为生成规约和证明节点的实际成本，$q_e$ 为检查调用边兼容性的成本。一次无返工 run 的工程量为
\[
  W_0=\sum_{v\in V}(s_v+p_v)+\sum_{e\in E}q_e.
\]
若把各操作视为单位成本，则三条路线都是 $W_0=\Theta(n+m)$。路线差异不在首次成功遍历的渐近阶，而在：何时暴露矛盾、多少已完成结果会因 bug 失效、需要保存多大的假设前沿，以及共享节点是否重复处理。

root spec 给定只省去或显著降低一个 $s_r$，不会改变渐近阶；但它能约束搜索空间，并改变 bug 被发现的时间。

\subsection{一次 Bug 传播的最坏成本}

对 spec bug $v$，节点化向上传播成本上界为
\[
  W_S(v)=O\bigl(|U(v)|+|E[U(v)]|\bigr),
\]
加上这些节点中已完成 spec/proof 的重做成本。对 proof bug，向下搜索上界为
\[
  W_P(v)=O\bigl(|D(v)|+|E[D(v)]|\bigr).
\]
最坏情况下 $U(v)=V$ 或 $D(v)=V$，所以单次 bug wave 为 $O(n+m)$。

\subsection{多轮返工的通用上界}

若共有 $b$ 个相互独立的 bug/revision wave，并且每轮都可能覆盖全图，则调度和重验工程量为
\[
  W=O\bigl((b+1)(n+m)\bigr)
\]
（未计 solver 内部复杂度）。若 $b=\Theta(n)$，则 $W=O\bigl(n(n+m)\bigr)$。稀疏 CG（$m=O(n)$）为 $O(n^2)$；稠密 CG（$m=\Theta(n^2)$）为 $O(n^3)$。

\subsection{三条路线各自的最坏返工形态}

\paragraph{路线 I。}
caller proof 先于 callee proof。链长为 $n$ 时，最深 callee 的 spec/proof 问题可能使此前 $n-1$ 个 caller proof 失效；若每深入一层都发生一次需要向上修正的问题，重做量为
\[
  1+2+\cdots+n=\Theta(n^2).
\]
一般图上，$b$ 轮最坏为 $O(b(n+m))$。路线 I 的主要风险是\emph{speculative proof debt}：已经完成但仍依赖未验证合同的 caller proof 数量。

\paragraph{路线 II。}
callee proof 先完成，因此不会因尚未验证的 callee 直接废弃已完成 caller proof。在 root spec 未给定、摘要推断单调且一次收敛时，最接近 $\Theta(n+m)$。但若 root spec 给定且直到最后才发现底层摘要不足，则一次顶层不匹配可触发全图向下加强；$b$ 次加强仍为 $O(b(n+m))$。此外，自底向上生成的路径敏感摘要可能随分支组合指数增长，这属于合同表示大小而非遍历次数。

\paragraph{路线 III。}
先固定规约，再从叶到根证明。callee proof bug 会在对应 caller proof 尚未开始时暴露，因此通常最少浪费已完成 proof；但如果证明失败要求修改 callee spec，spec bug 仍会向上使祖先规约 stale。$b$ 轮最坏同样是 $O(b(n+m))$，只是常数项通常更偏向重新生成 spec，而不是重新做已经完成的 caller proof。

\subsection{按路径传播会导致指数爆炸}

上述多项式上界依赖“节点访问去重 + 版本化边依赖”。若实现按调用路径而不是节点闭包传播，具有 $k$ 层二选一菱形的 DAG 可以拥有 $2^k$ 条 root-to-leaf 路径，而节点数只有 $O(k)$。此时同一 bug 会被重复报告/重验
\[
  \Theta(2^k)=\Theta(2^{n/2})
\]
次。工程实现必须：
\begin{itemize}
  \item 以 $(v,\mathsf{version})$ 为访问键；
  \item 对共享 callee 汇总入边需求；
  \item 每个 bug wave 对一个节点最多入队一次；
  \item 在 SCC 层面而不是递归路径层面提交结果。
\end{itemize}

\subsection{SCC 与候选规约搜索}

若 SCC $K$ 含 $r$ 个节点和 $e_K$ 条内部边，一次固定候选合同后的联合检查仍是 $O(r+e_K)$ 个调度操作；但找到合同不动点可能需要 $h$ 轮：
\[
  O\bigl(h(r+e_K)\bigr).
\]
若每个节点有 $k$ 个互不约束的合同候选，朴素联合枚举可达 $k^r$。Houdini 式候选删减通过反复调用 checker 删除不可证候选\cite{Flanagan2001Houdini}；ICE/CEGIS 则通过反例迭代学习不变量\cite{Garg2014ICE}。它们改变候选搜索策略，但不消除一般程序验证的不可判定性。因此必须区分：
\begin{enumerate}
  \item \textbf{CG 调度工程量}：在正确去重下按 wave 为 $O(n+m)$；
  \item \textbf{annotation 搜索工程量}：可能是多项式、指数或不收敛；
  \item \textbf{单个 VC 的证明复杂度}：取决于逻辑和 solver，一般不存在统一有限最坏上界。
\end{enumerate}

\section{路线比较与工程选择}

\begin{longtable}{>{\raggedright\arraybackslash}p{0.16\textwidth}
                  >{\raggedright\arraybackslash}p{0.25\textwidth}
                  >{\raggedright\arraybackslash}p{0.25\textwidth}
                  >{\raggedright\arraybackslash}p{0.24\textwidth}}
\toprule
路线 & 优点 & 主要风险 & 更适合的场景\\
\midrule
\endfirsthead
\toprule
路线 & 优点 & 主要风险 & 更适合的场景\\
\midrule
\endhead
I：caller proof 先行
& 很早验证顶层设计在假定接口下是否闭合；caller 驱动接口直接
& callee 合同不可实现时，祖先 proof 大面积失效
& root 需求最重要、callee 合同可快速稳定或来自可信团队\\
\addlinespace
II：bottom-up
& 每个 caller 只依赖已证明 callee；底层复用和增量性好
& 给定 root 规格时可能很晚才发现实现摘要不足；实现导出的 spec 未必表达意图
& 底层库先建设、历史代码摘要推断、root spec 未给定\\
\addlinespace
III：spec-down/proof-up
& 在 proof 前统一处理需求；callee 失败发生在 caller proof 前
& spec pass 可能生成之后证明不需要或不可实现的合同；仍需回流
& 高可信模块化验证、明确 root 需求、共享 callee 较多\\
\bottomrule
\end{longtable}

三条路线都可以合法；它们的核心差异是两个偏序：
\[
  \prec_S\;\text{（spec 顺序）},\qquad
  \prec_P\;\text{（proof 顺序）}.
\]
路线 I 允许 caller 的 $\ProofOp$ 先于 callee；路线 II 的 $\SpecOp$ 和 $\ProofOp$ 都是 bottom-up；路线 III 的 $\SpecOp$ top-down、$\ProofOp$ bottom-up。

\section{对 QCP 编排的直接含义}

若 QCP 要显式支持三条路线，controller 不应只记录 phase，还应为每个节点记录
\[
  (\mathsf{spec\ status},\mathsf{spec\ version},
    \mathsf{proof\ status},\mathsf{proof\ dependencies},
    \mathsf{bug\ state}).
\]
并提供：
\begin{enumerate}
  \item CG/SCC 构造和稳定节点 ID；
  \item $S/P$ 双访问序列生成器；
  \item 路线 I 的“条件完成 proof”状态，直到 callee proof 全部 discharge；
  \item 路线 II 的 bottom-up summary/spec 接口；
  \item 路线 III 的 top-down demand aggregation 和 reverse proof schedule；
  \item spec bug 的 reverse-reachability invalidation；
  \item proof bug 的 forward-reachability diagnostic search；
  \item 版本化、节点去重和 SCC 联合提交，避免按路径指数重复。
\end{enumerate}

现有 QCP 的单 case 流程更接近路线 III 的局部版本：先 annotation，再生成 VC，最后 proving；但只有把函数与循环组织成上述 CG 节点、明确 $S/P$ 两次访问及 bug 传播，才能讨论跨函数的三条完整宏观路线。

\section{结论}

对任意调用图，先缩并 SCC 得到 DAG，再选择三种调度之一：
\[
\begin{array}{ll}
\text{I}   & \text{spec-down，caller proof 可先行，callee proof 延后 discharge};\\
\text{II}  & \text{callee spec/proof 全部 bottom-up};\\
\text{III} & \text{spec top-down，proof bottom-up。}
\end{array}
\]
每个 function/loop 节点恰有一次当前版本的 spec 主访问和一次 proof 主访问；修复产生新版本和新 wave，而不是把同一旧版本沿每条路径重复遍历。无 bug 时三条路线的 CG 调度均为 $\Theta(n+m)$。单个全图传播 bug 为 $O(n+m)$；$b$ 轮为 $O(b(n+m))$；当 $b=\Theta(n)$ 时，稀疏图最坏 $O(n^2)$，稠密图最坏 $O(n^3)$。若不按节点去重而按路径传播，则即使 DAG 也可能指数爆炸。annotation 候选搜索和 VC 求解还有独立的指数或不可判定复杂度，不能与调用图遍历成本混为一谈。

\begin{thebibliography}{99}

\bibitem{Leino2010Dafny}
K. Rustan M. Leino.
\newblock Dafny: An Automatic Program Verifier for Functional Correctness.
\newblock In \emph{LPAR-16}, LNCS 6355, pp. 348--370, 2010.
\newblock DOI: \href{https://doi.org/10.1007/978-3-642-17511-4_20}{10.1007/978-3-642-17511-4\_20}.

\bibitem{Calcagno2009BiAbduction}
Cristiano Calcagno, Dino Distefano, Peter W. O'Hearn, and Hongseok Yang.
\newblock Compositional Shape Analysis by Means of Bi-Abduction.
\newblock In \emph{POPL 2009}, pp. 289--300, 2009.
\newblock DOI: \href{https://doi.org/10.1145/1480881.1480917}{10.1145/1480881.1480917}.

\bibitem{Lahiri2009IntraModule}
Shuvendu K. Lahiri, Shaz Qadeer, Juan P. Galeotti, Jan W. Voung, and Thomas Wies.
\newblock Intra-Module Inference.
\newblock In \emph{CAV 2009}, LNCS 5643, pp. 493--508, 2009.
\newblock DOI: \href{https://doi.org/10.1007/978-3-642-02658-4_37}{10.1007/978-3-642-02658-4\_37}.

\bibitem{Flanagan2001Houdini}
Cormac Flanagan and K. Rustan M. Leino.
\newblock Houdini, an Annotation Assistant for ESC/Java.
\newblock In \emph{FME 2001}, LNCS 2021, pp. 500--517, 2001.
\newblock DOI: \href{https://doi.org/10.1007/3-540-45251-6_29}{10.1007/3-540-45251-6\_29}.

\bibitem{Garg2014ICE}
Pranav Garg, Christof L\"oding, P. Madhusudan, and Daniel Neider.
\newblock ICE: A Robust Framework for Learning Invariants.
\newblock In \emph{CAV 2014}, LNCS 8559, pp. 69--87, 2014.
\newblock DOI: \href{https://doi.org/10.1007/978-3-319-08867-9_5}{10.1007/978-3-319-08867-9\_5}.

\end{thebibliography}

\end{document}
